\documentclass{article}
\usepackage{/globals/tables}
\usepackage{/globals/anki}
\ankiprefix{innovation}
\ankitopic{Innovation Adoption}
\ankishow
\begin{document}
\section{Adoption von Innovation} 
Die Adoption neuer Technologien kann durch die Diffusionstheorie erklärt werden. Diese Teilt die NutzerInnen in fünf Gruppen ein; die \emph{Innovators}, die \emph{Early Adopters}, die \emph{Early Majority}, die \emph{Late Majority} und die \emph{Laggards} (NachzüglerInnen).
\begin{ztable}{lXr}{Name & Gruppe & Größe}
 Innovator & Junge, technologieintressierte Personen; FreundInnen von Silicon Valley; Startups & 2.5\% \\
 Early Adopters & Personen die kein Problem damit haben BetatesterInnen zu sein & 13.5\% \\
 Early Majority & Personen die eher fertige Software wollen & 34\% \\
 Late Majority & Sozialer Druck die Technologie zu nutzen & 34\% \\
 Laggards & Ältere, technologieferne Personen & 16\% \\
\end{ztable}
Je nach dem wo auf dieser Skala eine Person (in Bezug auf eine bestimmte Technologie) befindet kann sie davon überzeugt werden diese zu nutzen. Die Innovators zum Beispiel müssen garnicht überzeugt werden, sie suchen nach dieser von sich aus, während die Late Majority nur mit sozialem Druck von Vertrauenspersonen dazu überredet werden kann.

\ankinote{gruppen}{Gruppen}{Innovators, Early Adopters, Early Majority, Late Majority, Laggards}
\ankinote{size-innovator}{Innovator Größe}{1/40}
\ankinote{size-earlyadopters}{Early Adopters Größe}{1/6 - 1/40}
\ankinote{size-earlymajority}{Early Majority Größe}{1/3}
\ankinote{size-latemajority}{Late Majority Größe}{1/3}
\ankinote{size-laggards}{Laggards Größe}{1/6}
\end{document}